\section{Οργανωτής GA (Orchestrator)}
\label{sec:impl_orchestrator}

Ο Orchestrator συνιστά τον οδηγό του γενετικού αλγορίθμου για κάθε νήσο. Υλοποιεί στρατηγική επιλογής $(\mu + \lambda)$ με ελιτισμό, διασφαλίζοντας ότι οι καλύτερες λύσεις διατηρούνται αναλλοίωτες μεταξύ των γενεών.

\subsection{Κύκλος Γενετικού Αλγορίθμου $(\mu+\lambda)$}
\label{subsec:impl_ga_loop}

Ο κύκλος εξέλιξης εκτελείται ως ακολούθως: αρχικά δημιουργείται ένας πληθυσμός $\mu$ ατόμων με τυχαία αρχικοποίηση εντός των ορίων των γονιδίων. Σε κάθε γενεά, εφαρμόζονται τελεστές τουρνουά επιλογής, ομοιόμορφης διασταύρωσης (uniform crossover) και Γκαουσιανής μεταλλαγής (Gaussian mutation) για τη δημιουργία $\lambda$ απογόνων. Μετά την αξιολόγηση — μέσω του αγωγού $\epsilon$-bypass — οι $\mu + \lambda$ συνολικοί υποψήφιοι ταξινομούνται κατά φθίνουσα καταλληλότητα και διατηρούνται μόνο οι κορυφαίοι $\mu$, εξασφαλίζοντας τον ελιτισμό. Η μεταλλαγή ακολουθεί κανονική κατανομή $\mathcal{N}(0, \sigma_m^2)$ με ρυθμό μεταλλαγής $p_m = 1/d$ ανά γονίδιο, όπου $d = 10$ η διάσταση του διανύσματος σχεδιασμού. Ο βρόχος επαναλαμβάνεται για $G_{\max}$ γενεές ή έως ότου το κριτήριο σύγκλισης — στασιμότητα της καλύτερης τιμής καταλληλότητας για $\Delta G$ συνεχόμενες γενεές — ικανοποιηθεί.

\subsection{Αποθήκη Γονιδίων GeneStore}
\label{subsec:impl_genestore}

Η τοπική αποθήκη ατόμων υλοποιείται ως \texttt{BTreeMap<HilbertKey, Individual>}, όπου \texttt{HilbertKey} είναι ένας 64-bit ακέραιος που προκύπτει από την κωδικοποίηση Hilbert του διανύσματος γονιδίων~\cite{skilling2004hilbert}. Η επιλογή BTreeMap αντί HashMap εξασφαλίζει $O(\log n)$ εισαγωγή και αναζήτηση, αλλά κυρίως ενεργοποιεί αποδοτικά ερωτήματα εύρους (range queries) που αξιοποιούνται κατά τη δειγματοληψία γειτονικών ατόμων για διασταύρωση. Κάθε \texttt{Individual} αποθηκεύει το διάνυσμα γονιδίων, την τιμή καταλληλότητας, τον αριθμό γενεάς δημιουργίας και ένδειξη του επιπέδου αξιολόγησης (Tier 1/2/3) που χρησιμοποιήθηκε.

\subsection{Αυτόματη Εκκαθάριση Γενεών: GenerationEvictor}
\label{subsec:impl_evictor}

Η συνεχής εισροή νέων ατόμων κινδυνεύει να εξαντλήσει τη μνήμη RAM σε μακρόχρονες εκτελέσεις. Ο \texttt{GenerationEvictor} εκτελεί σάρωση εκκαθάρισης (eviction scan) ανά $G_{\text{evict}}$ γενεές: αφαιρεί από το \texttt{GeneStore} όλα τα άτομα με αριθμό γενεάς $g < g_{\text{current}} - W_{\text{ttl}}$, όπου $W_{\text{ttl}}$ το παράθυρο TTL (Time-To-Live) σε γενεές. Η εκκαθάριση εκτελείται σε ξεχωριστό νήμα Tokio (\texttt{tokio::spawn}) ώστε να μην επιβαρύνει τον κύριο βρόχο εξέλιξης. Πριν από την αφαίρεση, τα κορυφαία $K_{\text{elite}}$ άτομα της τρέχουσας γενεάς αποτυπώνονται σε snapshot και αποστέλλονται στο LEAD DHT για μόνιμη αποθήκευση, διατηρώντας ιστορικό βέλτιστων λύσεων.

\subsection{Διακόπτης Κυκλώματος (Circuit Breaker)}
\label{subsec:impl_circuit_breaker}

Για την ανθεκτική αντιμετώπιση προσωρινών αποτυχιών δικτύου, κάθε εξερχόμενη gRPC κλήση προς εργάτες ή surrogate διέρχεται από έναν διακόπτη κυκλώματος (circuit breaker) με τρεις καταστάσεις:
\begin{itemize}
    \item \textbf{Closed (κανονική λειτουργία):} Κλήσεις διαβιβάζονται κανονικά. Εάν ο αριθμός αποτυχιών σε κυλιόμενο παράθυρο $W_{\text{cb}}$ υπερβεί το κατώφλι $F_{\text{thresh}}$, ο διακόπτης μεταβαίνει σε Open.
    \item \textbf{Open (διακοπή):} Κλήσεις απορρίπτονται άμεσα χωρίς δοκιμή αποστολής. Μετά από χρονικό διάστημα ανάκαμψης $T_{\text{rec}}$, ο διακόπτης μεταβαίνει σε Half-Open.
    \item \textbf{Half-Open (δοκιμαστική λειτουργία):} Επιτρέπεται μία δοκιμαστική κλήση. Εάν επιτύχει, ο διακόπτης επανέρχεται σε Closed· εάν αποτύχει, επιστρέφει σε Open.
\end{itemize}

\subsection{Μετανάστευση σε Δακτύλιο μέσω gRPC}
\label{subsec:impl_migration}

Ανά $T_{\text{mig}}$ γενεές, ο Orchestrator σειριοποιεί τα κορυφαία $K_{\text{mig}}$ άτομα και τα αποστέλλει ασύγχρονα στον διάδοχο κόμβο του δακτυλίου μέσω κλήσης \texttt{MigrateIndividuals(MigrationRequest)}. Η σειριοποίηση γίνεται με Protocol Buffers για ελάχιστο μέγεθος μεταφοράς. Η αποστολή είναι fire-and-forget: ο Orchestrator δεν αναμένει επιβεβαίωση παραλαβής, αποφεύγοντας μπλοκάρισμα του κύριου εξελικτικού βρόχου. Σε περίπτωση αποτυχίας, ο διακόπτης κυκλώματος καταγράφει το συμβάν και η μετανάστευση παραλείπεται για τις επόμενες $T_{\text{mig}}$ γενεές.

\subsection{Επανάληψη με Ποινή Αποθεματικής Λύσης}
\label{subsec:impl_retry_fallback}

Όταν όλοι οι διαθέσιμοι εργάτες VLM είναι κατειλημμένοι και η ουρά αναμονής υπερβεί το όριο $Q_{\max}$, ο Orchestrator εφαρμόζει στρατηγική fallback: το άτομο επιστρέφεται με τιμή καταλληλότητας ποινής $F_{\text{penalty}} = -\infty$, σηματοδοτώντας ότι δεν αξιολογήθηκε. Στην επόμενη γενεά, το άτομο αυτό επανεισάγεται στην ουρά αξιολόγησης με αυξημένη προτεραιότητα. Αυτός ο μηχανισμός αποτρέπει τη σιωπηλή απόρριψη ατόμων και διατηρεί την ποικιλομορφία του πληθυσμού ακόμα και υπό συνθήκες κορεσμού των πόρων.

